\chapter{第二集 \quad 摇篮}

\section{正文}

天地玄黄，宇宙洪荒，日月盈仄，辰宿列张。我们的中国正从远古混沌中渐渐清晰。6,500万年前青藏高原隆起，黄土高原、云贵高原、内蒙古高原，抬升形成了中国三级阶梯式的基本地理格局。自此，大江东去，日月西沉，疾风地北，烟雨江南，是我们所熟悉的中国。距今约 200万年前，直立人出现在东亚大陆东方古人类繁衍生息持续演化的故事拉开帷幕。

云南元谋人、陕西蓝田人、北京人、广东马坝人、山顶洞人，我们只能依据发现地点给这些远古的祖先命名。他们离开我们太过久远，久远的骨骼已经石化，还执拗地保留着中国古人类独有的特征。头骨正中的矢状脊，突出的面部，高而前凸的颧骨、阔鼻、铲形上门齿、下颌的圆枕．.

北京房山周口店遗址的20多个地点都发现了古人类生活过的遗迹。在位于第27 地点的田园洞地层里，出土了距今约4万年的人类遗骸化石。其中包括下颌骨、脊椎骨、股骨跟骨等体骨的大部分骨骼，属于同一位中年男性个体。

古人类遗骸或遗迹中残存着极其微量的 DNA片段，让我们可以直接研究过去人群的遗传信息。正是从田园洞人的腿骨上，古遗传学家成功捕获到测序较完整的基因组序列，与现代人基因特征完全符合。这是中国第一例人类古基因组，也是目前为止获得的东亚最古老的人类基因组。遗传意义上的亚洲人的祖先，在距今4万年前已经出现，他们与世界各地现代人的出现基本同时。

然而，挑战在3万年前后降临，严酷的冰期席卷全球，随着气温不断下降，巨大的冰盖覆盖了北美、欧洲和亚洲北部的大部分地区。

曾经水草丰美，猛犸成群的草场突然变成了一望无垠的冰原和荒漠；曾经湿热瘟瘴的丛林，却可能变成了气候宜人的河谷。适宜人类生存的区域不断偏移和缩小，风雪交加的世界里，田园洞人的身影渐渐模糊．...

约两万年前，江西万年县大源盆地的山间，一群人正在黄昏的光线里寻找着入夜后的栖息之所。眼尖的少年看见小河山的脚下，因为河床的缩小和后退，新露出了一处山洞。山洞内外地势平常，离水源不远也不近，是个十分理想的住所。于是人们对洞内的地面稍作平整后，生起一堆火，在这里度过了大雨滂沱的寒冷夜晚，这是被称为末次冰期最盛期的时代。环境发生着急剧的变化，熟悉的食物逐渐消失，现代人面临着自出现以来最严峻的生存挑战。

初升的日光冷清而凛冽，仙人洞族人已早早开始了一天的忙碌。老人敲击燧石砸出石片，锋利的刃部足以割开兽皮；女人用骨针缝补兽皮，制成过冬之衣；族人们收获野生禾草，用骨鱼杈捕鱼，也捕捞河蚌和田螺；孩子们则在原野和林间操练着捕猎的技能。

河岸边的少年就地取用夹杂着石英砂砾的泥土，摻水做成泥条一圈圈盘筑成園底的深腹罐。他在罐内外不停的拍打，使得泥条间的粘合更加规整、紧密。天色越来越明亮，少年将做好的罐子晾在一边，再将另一些早已阴干的陶坯放进露天的火堆中，渐渐的泥土变得干燥而坚硬。

陶器是数百万年的演化以来，人类创造出的第一种人工材料。这是目前已知世界上最早的陶容器﹣﹣仙人洞出土的陶器残片。外表面仍留存有烟食与火烧的痕迹，这显示他们应该是炊具。从此在生食和烧烤之外，人类还可以通过粒食和蒸煮的方式加工食物，从富含淀粉的食材和肉类中获取更多的能量。

外出的人们带着一天的收获归来，山洞里再度热闹起来，有人用锋利的刮削器肢解了一头斑鹿，放在火堆上炙烤。孩子们围陶罐而坐，煮熟的食物香气扑鼻，有人敲开鹿骨将骨髓和另一些软烂的食物留给老人。

天气越来越冷，食物的获取并不稳定，不是每天都能像这样饱餐。因此所有的食物都要物尽其用。在蒸腾的热气中，山洞里的火光明明灭灭仿佛悬挂在远古文明之路上的一盏夜灯，远在数百万年前当古猿人开始直立行走那一刻，他们看到了更远的世界，却也暴露出柔软的腹部。此后，漫长的狩猎采集生涯里，更需要彼此守护。

广东英德青塘 13,500年前的拂晓，一位豆蔻年华的少女刚刚离世，同伴们将她抬出日常活动的洞穴，穿过一片疏林来到另一个狭长幽深的洞中。他们选了一处凹坑，将少女摆出蹲踞的姿势放进去，在她背后垫了六块石灰岩角砾，把她生前常常使用的一枚骨针也放了进去。

这一次仍然生存的人们心头涌起了悲哀和思念，他们不能再像其他的动物一样任由死去的同伴在自然中慢慢分解，他们想用这样一种"仪式"来告别曾经一同生活过的人，让她用一种特别的方式沉睡，就像婴儿回到了母体，就像活着的人和死去的人之间永远有一种联结。

一万年前随着全新世的到来，气候逐渐变暖，人类的生存条件好转，上山文化的先民自洞穴走向旷野，在浙江的金衢盆地停留下来，沿着钱塘江的支流营建房屋和村落，过上了定居的生活。

村庄里有经验的前辈们，春天在这里试着播下种子，随着雨季后水位下降，裸露的河岸边如今结出了饱满的稻穗。少年收割稻穗，更年幼的孩子也来帮忙，用手轻拍还没完全驯化的谷粒就纷纷落下了。母亲们用磨盘、磨棒，碾磨橡子、稗子、菱角、块根放到陶制的大口盆中，加热搅动后就能煨成糜烂的粥。

制陶匠人把散落的稻壳、稻叶和稻杆掺在陶土中，可以增加陶土的延展性，提高耐热性能，避免陶器在烧制过程中发生破裂。这些掺杂着稻壳和稻叶的陶片，成为植物考古学家研究的对象。

通过对小穗轴基盘形态的观察，可以确认有些稻穗不是自然脱落的，已有人工驯化迹象。这是世界上最早的碳化稻种子，距今9,000 多年发现于浙江上山遗址 461号灰坑的填土中，它可能是古人在炊煮过程中，偶然掉落在火塘边的。虽然有些残缺，但它仍然保留了早期稻作农业的珍贵信息。

早在更久远的时代，人们就已经认识到野生稻种子能用来果腹，或许他们无意间发现生长于河岸边的这些稻属植物在第二年还能发芽抽穗，于是尝试着最初的种植。尽管作物的完全驯化需要经历数千年的演变过程，又或许稻米在当时人们的食谱中仍非主导，但农耕的出现却永久的改变了历史的走向。

在上山文化延续的近两千年间，炊烟袅袅，人口蕃息……在浦江上山，人们修建了 10 余米长的干栏式房屋，在义乌桥头聚落三面环壕，一面临河，其间发现了房屋、墓葬和窖穴遗迹，反映出丰富的定居生活内容。

安定祥和的生活里，长出了从容有度的生命，比之一万年前就连陶器的面貌也焕然一新，多种形态的盆、罐、钵、杯、盘是为了应对更精细的生活需求。

以白彩施于红底之上的太阳纹，对顶三角纹形似八卦卦象的短线组合纹等，则投射了更深邃的精神世界和更丰富的审美情趣。

这件义乌桥头遗址出土的陶壶，虽历经8,500年岁月斑驳，仍难掩动人魅力，造型简约而优雅，线条流畅而柔和，陶色明丽但不炫目，静静凝聚成历史那头的一抹红。

与此同时，定居社会与农业革命的强音也在北方奏响。河北西北部的张家口尚义四台遗址发现成排半地穴式房址，测年数据达到距今1万年左右，是北方地区最早的明确的定居聚落。遗落在地面上的打制刮削器，工艺成熟的细石器以及大量的动物骨骼，表明食物种类变得更加丰富，但狩猎采集仍然是获取食物的重要手段。磨盘、磨棒和陶器的出现则显示食材的精细加工和烹制已开始普及，定居生活方式的出现成为旱作农业在北方地区起源发展的重要条件。

从目前的考古证据来看，北方旱作农业最早的发生区并非中原，而是更靠北的燕山南北地带。粟和黍都在中国传统的"五谷"之列。北京西郊的东胡林遗址中，发现了迄今所见年代最早的碳化粟粒。粟就是小米，田间地头随处可见的狗尾草，是粟的野生祖本。狗尾草的籽粒狭长，经过人类的驯化，粟粒的形状已经接近球形。

阴山山脉以北科尔沁沙地的边缘地区，兴隆洼的先民们则已经开启了以黍为主要农作物的村落生活。由于它耐寒、耐旱的特点，黍在8000年前的整个北方地区，成为了最重要的旱作农业品种。

今天每一种出现在我们餐桌上的食物，都曾见证了远古先民们的饥寒与艰辛，他们播下稻种、粟种，从此不再被动等待自然的赐予，而是主动参与到万物生长的进程中。中华文明以农为本的基因早已渗透在我们的思维和血液里。

8000年前兴隆洼猎人的队伍，穿行在西辽河流域的平原丘陵间。尽管这里的人们已经过上了以种植粟黍为主的农业生活，但面对将到来的漫长冬季，通过狩猎来储备食物依然不可或缺。

英武的首领是聚落中最优秀的猎手，正带领着族人向丛林深处行进，人们身背长弓，腰悬兽皮箭囊。首领示意猎人们在草丛中隐蔽下来，他侧耳倾听又仰面观察树叶的摆动，前方忽然闪现一只野猪的身影。首领摘下长弓，抽出羽箭。

这令人惊叹的考古现场正是兴隆洼人曾经的家，它的面积达3万平方米，围沟环绕中 180余座半地穴式房屋成排分布，从小型流动的狩猎采集人群，发展为规划有度的大型村落社会。兴隆洼人经历了巨大的环境、资源和社会的挑战。

当年在聚落中心最大的房屋里，首领正在接待远方来的贵客，众人围坐在火塘边相谈正欢。首领拿出一柄刃部装嵌有细石叶的骨埂石刃刀，这是他爱用的狩猎工具。

8,000年后考古学家在这间房屋的东北部地面下，发现了一座墓葬。墓主人是一位50多岁的男性猎人首领。把少数生前或许具有特殊地位的人埋在居室内是兴隆洼文化的特殊葬俗。墓主右侧葬有两只整猪，一雌一雄均呈仰卧状，似乎是祈求猎物丰盛之意。在他的耳部发现了一对玉玦，根据使用痕迹来看应该是墓主人非常珍爱并日常佩戴的耳饰。

玉是美石，经过琢磨、抛光后温润而莹泽，后来的中国发展出了世界上独一无二的玉文化，此为重要先声。

在农业起源的初期，中华大地已经呈现出多元的发展景象，南北方形成了各具特色的农业模式。淮河流域则表现出南北交汇，错落相间的文化格局。

距今 9,000 年到 7800 年前生活在河南舞阳县贾湖遗址的先民，因附近有湖泊湿地，鸟兽成群，便以稻作、渔猎为生。少年蹲在莎草丛边的石上，用湖水仔细的清洗着一段丹顶鹤翅膀上的尺骨，身后的父亲接过尺骨，用手指认真的比度后，刻下几道细线交给少年，少年根据标记用燧石尖刃器小心地在尺骨上钻下了7个孔。

骨笛的制作涉及到复杂的音律和数学知识，在氏族里只有少数人能传承古笛制作和吹奏的技艺。一旦制成，骨笛也往往永远伴随着主人生死不离。

贾湖村落中不同家族的房屋各自成组，错落环绕在中心广场周围，房屋之旁还有制作陶器的窑场，储藏食物的窖穴，死去的人们则安睡在离房屋不远的地下。

墓葬是我们解读过去的锁钥。在贾湖二期的墓葬里344号墓的墓主为男性，无头，代之以叉形器和一组龟甲，其中一件龟甲上还有意义不明的刻符，左臂旁放着一件骨饰﹣﹣两只骨笛，下肢处则放置一些渔猎工具。推测墓主人可能是身具多种才能的智者，担任着氏族的族长兼巫师。

尽管比起数万年前生活已经好过了一些，但人们种植水稻、饲养家猪的规模相当有限，很大程度上还需要依靠渔猎采集获取食物，冬春季节分外难捱。

于是这夜有一场隆重的仪式，氏族成员们集中在村落中央的广场，族长和女祭司肃立其间，高亢的骨笛之声划破夜空。年轻人将酒倒入钵中，用大米、蜂蜜和山楂酿成的酒，甜香沁脾，族长一饮而尽，手持刻符石柄的女祭司高举起双手。成年族人的小腿上绑着装有石子的龟甲响铃，随着舞蹈的节奏怦然作响，少女们的颈下和腰间挂着一串串鸟肢骨管横截而成的小骨环，舞姿轻灵。通过这样的仪式，族人们祈求神灵庇佑；祈求年岁丰穰；祈求氏族人口繁衍；祈求春来水患不扰．...

聚落之外，少年静坐河岸，月光如水洒满河滩，少年抽出那支崭新的七孔笛放在唇边。祖先常在，神灵常在，四时循环，万物化生，万籁俱寂中少年感受到大音希声之妙境，风声、水声、芦苇声、鹤鸣声共作，涌到唇间，送入笛身。他越奏越入神，手指按向骨笛的第7孔。广场上忙碌的人群都为此声惊喜，静止不动，静听余音在星空中盘旋．....

时光荏苒，昔日的少年成为了新的族长。他已经拥有了自己独立制作的骨笛，但仍细心保留着已经故去的父亲当年教他制作的第一支骨笛。那只骨笛后来被失手摔成三段，他便在断裂处钻上几对缀合孔，用线细细缠裹，音色仍旧优美如初。

杭州湾畔的跨湖桥遗址，湖岸边的独木舟被固定在桩架设施中，奖置于船边，这样一停就是8,000年，似乎仍在等待起航。独木舟取材自一整棵的马尾松，制作时先用火烘烤木材表面，然后再用坚硬的石器一点点地木为舟，至今残长犹有5.6米。船旁发现了一些编织物，或许曾是一张船帆，在等待扬帆起航的日子里，它被海浪吞没。那是来自远古的一个信号，面对风浪，面朝大海，面向远方已经站立起我们的少年中国。

8,300 年前古宁波湾的海平面上升，海水灌入四明山北麓，余姚井头山遗址距离现在的海岸约20多公里，当年却是依山傍海。晨雾散去，岸边高地上露出一座村落，用粗木搭建起框架，地板架空在地面上，以竹篾为墙，茅草为顶，这样的干栏式房屋正适合温暖潮湿的海边环境。

井头山人享用了海鲜大餐后，就把鱼骨和蛎、蚝、螺的壳丢弃在村落旁。日复一日，年复一年，堆积如丘，这就是考古学中所称的"贝丘遗址"。

晨光里有人在海岸边踏浪，拾起还来不及随潮水褪去的贝类；有人双手翻飞用芦苇编织着背篓和渔罩；有人忙着收拾晾晒好的肉脯和鱼干，将橡子、麻栋果和核桃一起放进储藏坑，而拥有高超技艺的木匠们，自如的制作着各类木器。

我们所以为的历史往往与权力、战争和王朝更替相关，其实真正支撑人类数千年文明与温暖绵延的，是那些昼出夜伏，炊烟袅袅的日常生活。

石刀刮削木材的声音打破宁静，父亲在修整刚制作好的木桨，他校正杆部是否笔直，桨部的角度是否合适。父亲深知桨在航海中的重要性，在陆地上沿着嶙峋的海岸寻找合适的滩涂，费时且艰难，人们想依靠舟船去更远的地方。少女抬头向东望去，大海是一片灰白而忧郁的原野。

井头山遗址的库房中，这只木桨保存的完好程度令人吃惊，它好像刚刚完成了一次出海。发掘者在10米深的地下找到了这处遗址，这里后来曾被海淹没，覆盖着厚达8米的淤泥，淤泥容易流动坍方，因此在发掘前需要预建钢结构围护的发掘基坑。如果不是考古学家的远见和坚持，我们不会知道这海相沉积之下，曾也是一群人的家。

人与海进退之间，相伴相生。千年之后海平面趋于稳定，陆地重新露出，淡水溪河再次流淌于平原之上，成为人类宜居的家园。一叶竹筏在落日余晖中，摇进了青山绿水环抱的村庄。河姆渡同样的地理坐标，经历海岸线变化下的古环境演变，已然沧海桑田。成排的干栏式木屋错落于菱荷丛生的河湖溪塘间，低洼的平地上稻穗金黄，芦荻作雪飞，小桥流水，饭稻羹渔，颇有几分后世的江南模样。

河姆渡先民制作着翻土的骨耜，精耕细作成为这片土地的传统。与此同时人们对生活更多了分精致的追求，炊煮食物用釜，蒸食物用甑，煮水用盘，盛放食物用豆、盘、钵和盆。简单蒸煮食物以求果腹的时代已经过去，这些优雅美观的饮食器具的制作者似乎正于日常生活中感知着烹饪的愉悦。

这只圈足木碗外表还晕着朱红色的漆。时光已逝，考古学使它复活，那消失的生命的温度重返人间。蚕纹象牙杖端饰、黑陶朱纹钵，这些堪称原始艺术的物件都是先民智慧与信仰的结晶。

这件雕刻"双鸟朝阳纹"的象牙蝶形器，器身正面中央阴刻五环重圈纹，周围绕以火焰状光芒，两侧各有一只鸟昂首相望，又似在引吭啼鸣。

远古先民在漫长的采集渔猎和农业活动中，观察季节变化与万物生长，体验寒暑交替和昼夜节律，产生了对太阳的敬畏与向往。当他们把这种情感寄托在如鸟一般能通天达地的动物身上，神灵崇拜就应运而生。

一轮红日正冉冉升起，驱散了山间的溟漾晨雾，沅水在湘西的千丘万壑间蜿蜒穿行。远古的匠人以深红色的着色剂，用软笔在白陶簋形器的外底部绘出一个太阳纹的图像，如从远古幽冥中穿云破雾而来，瞬间光芒万丈。

这件亚腰白陶罐，肩部饰篦点凤鸟，雄健的勾喙，敏锐的环眼，张扬的双翅，展示着非凡的神采。它的身侧是蝶形兽面图像，似乎蕴含庄严的神力。这类刻画着神圣图像的白陶器是高庙文化的典型器物。在这件白陶罐的领部由戳印细篦点构成的双羽翅獠牙兽面纹，构图复杂，如同一张长满利齿的阔口，出獠牙，还伴有羽翼和羽饰。兽面两侧各有一座云梯萦纤的"高阙"，它是不是远古祭祀中的天梯形象呢？

7,000多年前，这些构思诡谲的通灵祭器犹如一道天光照亮了幽暗的史前丛林。它产生于远离江湖平原的河谷山川，那里充满了神秘的色彩和浪漫的情调。在洪水频发和瘴疫横行的年代里，疾病与死亡的威胁无处不在。为了获得神灵的恩赐与庇佑，先民奉献出最虔诚的信念和最精湛的技艺，以获得心灵的慰藉。

高庙文化并不是一个人丁兴旺、持续长久的考古学文化，却凭借一群狂热的天才"艺术家"和他们创造的神秘作品在历史上留下了独特的影响力。

随后以洞庭湖为中心的汤家岗文化继承了它的传统，在他们的白陶器上大量涌现各种复杂的几何纹样。有学者猜测，八角星纹代表一种宇宙观，它常见于圈足盘的底部，把这些器物倒扣放置正如天体的模型，中间的圆形是天，外侧的八角则象征大地的八方。无论如何，器具本身只是一种实物形式的载体，它真正要表达的是人们寄予这种载体之上的精神信仰。原始信仰驱散了许多萦绕在人们心中对未知世界的恐惧，也成为不同人群、不同聚落之间连接的纽带。

从汤家岗文化所在的长江流域向北越过绵延的秦岭，就到了地理意义上的北方。秦岭北麓发源的渭河汇入黄河之处，正是辽阔的"八百里秦川"也是丰沃的关中平原。

沪河东岸的半坡村落，一座半地穴房屋刚刚建造完成。屋顶新铺的茅草散发着清香，男人抬入盛满谷物的大陶瓮放在后墙边，又将陶盆、陶钵、陶罐、石斧、石铲等一一摆放好。进门的左右两侧各有一个矮矮的土台子，那便是床铺。女人给它铺上新编的草席和松软的兽皮，火塘中灶火正燃，褐色的夹砂陶罐架在灶上，外表布满烟苔，里面的小米粥已经煮熟，颜色金黄诱人。

距今 7,000年至6,000 年左右，新石器时代的社会得以显著发展。仰韶文化在黄河流域开启了属于自己的时代，其早期被命名为半坡类型，一个个围沟环绕的农业村落像点点繁星一般出现在黄土大地的川谷盆地间。

半坡村落的中心广场上燃起一堆篝火，一只整猪被对半剖开放在木架上烧烤，浓香四溢。新居收拾停当，庆贺乔迁的宴饮拉开序幕。聚落中的男子们在长长的草席上落座；氏族的长辈们被安排在首座；妇女们端出烹制好的食物，有小米饭、羹汤还有烤熟的猪肉。男人们抱起小口尖底瓶浑浊的米酒倒满一盆，他们用小杯舀酒向老者致敬，然后一饮而尽。女人和孩子们也开始享用丰盛的食物，整个聚落一片欢声笑语。

考古学家在半坡聚落的壕沟内发掘出45座房子，200 多个窖穴以及牲畜围栏遗迹。整个居住区以一条小沟分为南北两片，中间有道路相通，壕沟以北有公共墓地，以东则有公共窑场。每个房屋内居住着一个家庭，每个单元是一个扩展家庭，而整个聚落就是一个氏族。

临潼姜寨遗址更清楚地表现出聚落规划中对亲属关系的关注，圆形围沟内房屋有100多座，明显分成5个单元，围绕着近4,000平方米的中心广场分布。每单元都有大中小型房址，所有房门都朝向中心广场，在这样的空间格局中每一个居住者都时刻感受着个人与家族部落的血脉联系。

逝去的亲人按照亲属关系分组埋葬在聚落周围。考古学家发现一种特殊的埋葬方式对血缘的强调近乎极致。河南灵宝的城烟遗址发现多座多人合葬墓，均属仰韶文化早期。人骨成层安放，其中第94号墓内共有人骨19具。这些人骨是从之前不同时代的墓坑中被迁移而来，重新埋葬在一起，似乎更多在彰显群体的凝聚力和亲密关系。生则共居，死则同葬，血缘不会被死亡割断。

皓月当空，半坡聚落的仪式进入高潮，酒酣之际有人吹起陶员，悠悠咽咽在夜空回荡。少年被埙声打动，悄悄起身独自回到居室。那置于高处的彩陶盆是少年的牵挂，它是为聚落里逝去的婴儿准备的，盆的内壁用黑彩绘出人面鱼纹。

仰韶文化半坡类型的彩陶中最有代表性的象生图案便是鱼纹，有些是单独出现而较为写实的鱼纹；有些是以直线与弧线描绘，圆点、弧线和弧边三角穿插而显得活泼灵动的简化鱼纹。纹样格式除平展式外，还出现了回旋、跳跃等姿态。

半坡类型晚期，彩陶上的单独鱼纹采取了夸张变形的艺术处理，变成上下对称的式样，趋于几何化。鱼与人面相结合的形象是半坡彩陶独有的图案，如著名的人面鱼纹，人和寄寓又互相转借，意味着人和鱼是交融的共同体。

被人格化了的鱼类图像和各式鱼纹，可能具有半坡氏族保护神的性质，将早夭的婴儿装殓在陶瓮中，以瓮为棺，以盆为盖，埋入土中，这是一种叫瓮棺葬的习俗。人面鱼纹绘制在陶盆内壁，当它倒扣在瓮口的时候，画案只有其内的亡灵才能看得到。瓮棺上还特地留出一个圆孔，让孩子的灵魂可以自由出入。

百万年来无边的夜空曾令我们的先祖们畏惧，猛兽、寒冷、黑暗，总是一再侵袭的饥饿，随时突然降临的死亡使每一个夜晚都曾如此艰难和漫长。直到第一双拾取火种的手开始让光亮在寒夜驻守；第一双制作工具的手开始让万物可被改造利用；第一双撒下谷种的手开始为来年种下希望。

为了生存我们的祖先一次次走向远方，平原、山地或是丘陵、海岸，他们披荆斩棘，筚路蓝缕，聚族而居，生生不息。从此这片苍茫大地上夜空深邃，群星闪耀，不再惧怕黑夜的人终于能抬头仰望星河。


\section{思辨句式:我们所以为的***，然而，真正***，恰恰是***}

\subsection{原文}

我们所以为的历史往往与权力、战争和王朝更替相关，其实真正支撑人类数千年文明与温暖绵延的，是那些昼出夜伏，炊烟袅袅的日常生活。

\subsection{思辨结构解析}

提出大多数人认同的常识，进一步分析这背后令人吃惊的本质。（甚至相反）

\subsection{写 2024 新高考一卷}

\textbf{作文题：} 随着互联网的普及、人工智能的应用，越来越多的问题能很快得到答案。那么，我们的问题是否会越来越少？

\textbf{仿写：} 我们所以为的技术革新总是与算法、数据和机器进步相关，其实，真正在暗流涌动中引领文明破浪前行，赋予人类以跨越式解放与蜕变的，是那根孜孜不倦、探索无界的好奇心，是那些深邃不息、发现并挑战未知之难越的灵魂。


\subsection{写 2024 全国甲卷}

\textbf{作文题：} 每个人都要学习与他人相处。有时，我们为避免冲突而不愿表达自己的想法。其实，坦诚交流才有可能迎来真正的相遇。

\textbf{仿写：} 我们所以为的缓解冲突往往与保持沉默、息事宁人相关，其实，真正能够化解矛盾、消融分歧、弥合人心裂痕的，是那些坦诚相待、推心置腹的深入交流。在冲突面前选择沉默，只会致使矛盾在岁月的冲刷下愈发根深蒂固；选择坦诚交流，才能共同耕耘出理解的沃土，迎来真正的相遇。


\subsection{写 2024 北京卷}

\textbf{作文题：} 几千年来，古老的经典常读常新，杰出的思想常用常新，中华民族的伟大精神历久弥新……很多事物，在时间的淬炼中，愈显活力和价值。请以“历久弥新”为题目，写一篇议论文。

\textbf{仿写：} 我们往往将历史的磅礴归因于那些璀璨的经典、深邃的思想与崇高的精神，然而，真正赋予“历久弥新”以深刻内涵的，恰是那些经过时光雕琢，在人民群众日常琐碎中不断焕发生机与意蕴的生活细节。


\subsection{写 2024 天津卷}

\textbf{作文题：} 在纷纷的世界中，无论是个人、群体还是国家，都会面对别人对我们的定义。我们要认真对待“被定义”，明辨是非，去芜存真，为自己的提升助力；也要勇于通过“自定义”来塑造自我，彰显风华，用自己的方式前进。

\textbf{仿写：} 我们常常误以为他人的定义即是我们存在的全部意义，然而，真正决定我们本质与未来的，恰恰在于我们如何审慎对待外界的“定义”，明辨是非，去芜存真，以此作为自我提升的助燃剂；更在于我们是否有勇气进行“自我定义”，以独树一帜的姿态塑造并展露内在风华，坚定不移地沿着自己选择的道路砥砺前行。


\subsection{写 2024 上海卷}

\textbf{作文题：} 生活中，人们常用认可度判别事物，区分高下。请写一篇文章，谈谈你对“认可度”的认识和思考。

\textbf{仿写：} 我们惯于以浅薄的认可度作为评判事物高下的准绳，却忽略了其背后更深层的本质与意义。其实，真正剖析“认可度”的内涵，我们会发现，它也可能是盲从与主观臆断的结合体，而事物的真正价值，应超越外界的喧嚣认可，直抵其内核的深刻与独特。
